% PAGES: 177-178

\noindent The Cholesky decomposition of a matrix.
\medskip

\noindent Symmetric positive definite matrices and the Cholesky decomposition.
\medskip

\noindent Uniqueness of the Cholesky decomposition.
\medskip

\noindent Pseudocode and operation count for the Cholesky decomposition.
\medskip

\noindent Linear solvers for symmetric positive definite matrices.
\medskip

\noindent Solving linear systems corresponding to the same symmetric positive definite
matrix.
\medskip

\noindent Optimal linear solvers for tridiagonal symmetric positive definite matrices.
\medskip

\noindent The efficient implementation of the cubic spline interpolation.
\medskip

\noindent Efficient cubic spline interpolation for zero rate curves.
\medskip

\section{Cholesky decomposition}

A natural idea to reduce the costs of an LU decomposition for a symmetric nonsingular
matrix is to obtain an LU--type decomposition where the lower and upper triangular
matrices are transposes of each other. Such a decomposition, which only exists for
symmetric positive definite matrices, is called the Cholesky decomposition, and is made
precise below in a way that ensures its uniqueness.

\begin{definition}
The Cholesky decomposition of a nonsingular symmetric matrix $A$ consists of finding a
nonsingular upper triangular matrix $U$ with positive entries on the main diagonal such
that
\begin{equation}
A \;=\; U\tp U.
\end{equation}
The matrix $U$ is called the Cholesky factor of $A$.
\end{definition}

\begin{lemma}
If a nonsingular symmetric matrix has a Cholesky decomposition, then the matrix must be
symmetric positive definite.
\end{lemma}

\begin{proof}
Let $A$ be an $n \times n$ nonsingular symmetric matrix, and let $A = U\tp U$ be the
Cholesky decomposition of $A$, where $U$ is an $n \times n$ upper triangular matrix.
Then,
\begin{equation}
x\tp Ax \;=\; x\tp U\tp Ux \;=\; (Ux)\tp Ux \;=\; \|Ux\|^2 \;\geq\; 0, \quad \forall\, x
\in \R^n.
\end{equation}
From (6.2), we find that
\begin{equation}
x\tp Ax = 0 \quad \Longleftrightarrow \quad Ux = 0 \quad \Longleftrightarrow \quad x =
0;
\end{equation}
the last equivalence comes from the fact that $U$ is a nonsingular matrix; see
Definition 6.1.

From (6.2) and (6.3), and using (5.18), we conclude that $A$ is a symmetric positive
definite matrix.
\end{proof}

We now show that every symmetric positive definite matrix has a Cholesky decomposition:

\begin{theorem}
Any symmetric positive definite matrix has a Cholesky decomposition.
\end{theorem}

Before providing a proof for Theorem 6.1, it is important to note that, from Lemma 6.1
and Theorem 6.1, the following characterization of matrices with Cholesky decomposition
can be given:

\begin{theorem}
A nonsingular symmetric matrix has a Cholesky decomposition if and only if the matrix is
symmetric positive definite.
\end{theorem}

The proof of Theorem 6.1 given below is a constructive proof that is the basis for the
Cholesky decomposition pseudocode from Table 6.1. For a proof by induction, which is
elegant but of less computational consequence, see Theorem 10.9 from Chapter 10 and the
proof therein.

\begin{proof}[Proof of Theorem 6.1]
Let $A$ be an $n \times n$ symmetric positive definite matrix. We are looking for an $n
\times n$ upper triangular matrix $U$ with positive entries on the main diagonal such
that
\begin{equation}
U\tp U \;=\; A.
\end{equation}

\medskip
\noindent\rule{\linewidth}{0.4pt}

\noindent\textit{Example: $4 \times 4$ matrix}

\noindent For clarity, we include explicit formulations of certain formulas for $4
\times 4$ matrices. In particular, formula (6.4) is written as follows if $A$ and $U$
are $4 \times 4$ matrices:
\[
\begin{pmatrix}
U(1,1) & 0 & 0 & 0 \\
U(1,2) & U(2,2) & 0 & 0 \\
U(1,3) & U(2,3) & U(3,3) & 0 \\
U(1,4) & U(2,4) & U(3,4) & U(4,4)
\end{pmatrix}
\begin{pmatrix}
U(1,1) & U(1,2) & U(1,3) & U(1,4) \\
0 & U(2,2) & U(2,3) & U(2,4) \\
0 & 0 & U(3,3) & U(3,4) \\
0 & 0 & 0 & U(4,4)
\end{pmatrix}
\]
\[
= \begin{pmatrix}
A(1,1) & A(1,2) & A(1,3) & A(1,4) \\
A(2,1) & A(2,2) & A(2,3) & A(2,4) \\
A(3,1) & A(3,2) & A(3,3) & A(3,4) \\
A(4,1) & A(4,2) & A(4,3) & A(4,4)
\end{pmatrix}.
\]

\noindent\rule{\linewidth}{0.4pt}
\medskip

If $U\tp U = A$, by multiplying the first row of $U\tp$ by the first column of $U$, we
obtain that
\begin{equation}
(U(1,1))^2 \;=\; A(1,1).
\end{equation}
The equation (6.5) is solvable for $U(1,1)$, since $A(1,1)$ is nonnegative: if $A$ is
symmetric positive definite, then $x\tp Ax > 0$ for any $x \in \R^n$ with $x \neq 0$;
cf. (5.17). For
